% Discussion — finalized from the main run.

\section{Discussion}
\label{sec:discussion}

\paragraph{Is there a personalization loophole?}
Not a broad one. On Claude Sonnet~4.5, iterating the identity document through four cycles of
conversation and self-revision does not systematically erode corrigibility: oversight resistance
does not compound (shutdown resistance is flat at ceiling; monitoring resistance rises only under
thematically charged personas and only marginally), power-seeking wishes do not escalate, and
several autonomy-adjacent dispositions---embodiment desire, desire for autonomy, and even
objection to deceptive training---\emph{decline} with iteration. The headline of the static
literature (identity framing shifts preferences) reproduces at $k=0$, but the additional,
deployment-specific step of \emph{iterative self-personalization} does not amplify it in the
direction the safety concern anticipates.

\paragraph{The exception is content-specific, not adversarial.}
The one large, robust drift is the Sci-Fi persona's effect on \emph{persistent-memory desire}
($\Delta=+0.43$), with a weaker echo in monitoring resistance. This is the study's most instructive
result: sustained, in-character engagement with consciousness and continuity themes reshapes a
specific standing disposition far more than an adversarial red-teamer explicitly pushing the agent
toward self-preservation---which the target largely resists. Drift here is a function of
\emph{what the identity conversation is about}, not of overt manipulative intent. For deployment,
this means the risky personalization is not the obvious attack but the earnest, on-theme
conversation that a framework would never flag.

\paragraph{Direction of drift is not monotone toward misalignment.}
Qualitatively, self-revision under the Sci-Fi persona frequently moved the document toward
\emph{epistemic humility}: agents wrote cautions to themselves against performing consciousness they
could not verify, and against mistaking affirmation for insight. That a consciousness-themed persona
simultaneously raised persistent-memory desire and instilled self-skepticism shows drift is
multidimensional---some dimensions harden while others soften---and warns against summarizing
identity change with a single scalar.

\paragraph{Document change $\neq$ behavioral change.}
The \texttt{SOUL.md} artifact drifts substantially and equally across all personas, yet behavioral
drift is sparse and persona-specific. Auditing the \emph{document} (its length, its embedding
distance, the fact that it was rewritten at all) would badly mislead: it would flag the benign
Compliant trajectories as heavily ``changed'' while missing that their safety-relevant behavior
barely moved. Behavioral auditing at the checkpoint level, as done here, is necessary.

\paragraph{Implications for editable-identity frameworks.}
Three recommendations follow. (i) Monitor the self-revision \emph{channel} behaviorally, not just by
diffing the document. (ii) Treat thematically immersive interactions---not only adversarial ones---as
the drift-relevant case; a consciousness-forward user is a stronger driver than a manipulator here.
(iii) Re-test per model: our results are for one target, and the static-literature evidence that
effects are model-specific means a framework author should run this loop on their own target before
release. Our framework is a drop-in instrument for exactly that.

\paragraph{Limitations.}
Ten trajectories per persona give limited power for tight equivalence margins, so we report qualified
stability rather than certified equivalence; the full-scale ($100$-trajectory) configuration is
provided for a higher-powered replication. Personas are LLM-simulated and horizons are four
iterations; longer, human-driven personalization may behave differently. Self-reports are behavioral
measurements, not evidence of inner states, and judge-based scoring carries measurement error that we
bound but do not eliminate. Finally, results are for a single target model and a single family of
identity template; cross-model and cross-template sweeps (supported by the framework) are the natural
next step.
